%% 第3巻　関数とデータ　本文（第12回編集委員会）
%% 章は \bkinputchap{ch3-NN} で読む（まだない章は、組むときに飛ばして「BKPAGE missing」を出す）
%% 第3巻は巻の色が黄なので、部扉の絵の線は白でなく bkOn（濃い紺）でかく
\providecommand\bkinputchap[1]{\IfFileExists{#1.tex}{\input{#1}}{\clearpage\typeout{BKPAGE missing #1}}}
\bkinputchap{ch3-00}

%% ======================== 第1部 関数とグラフ ========================
\bkPartOpener{1}{\R{第}{だい}3\R{巻}{かん}\hspace{0.8em}\R{第}{だい}1\R{部}{ぶ}}{\R{関数}{かん|すう}とグラフ}%
  {172/{「どこ」を2つの\R{数}{かず}で\R{伝}{つた}えられる？},
   160/{\R{表}{ひょう}・グラフ・\R{式}{しき}は、\R{同}{おな}じことを\R{言}{い}っている？},
   148/{ふたりが\R{出会}{で|あ}うのは、グラフのどこ？}}%
  {% 座標の軸と直線（濃い紺の線）
   \begin{scope}[bkOn, line width=1.2pt, line cap=round, line join=round]
     \draw[-{Stealth[length=2.6mm]}] (118,112) -- (172,112);
     \draw[-{Stealth[length=2.6mm]}] (122,108) -- (122,152);
     \draw[line width=1.6pt] (122,116) -- (166,146);
     \draw[line width=1.6pt] (122,140) -- (168,120);
     \fill (143.5,130.65) circle (2.2);% 第13回（依頼者の指摘）：点を2本の直線の交点に
   \end{scope}}%
  {\R{変}{か}わる\R{量}{りょう}どうしの\R{関係}{かん|けい}を、\R{表}{ひょう}・グラフ・\R{式}{しき}の3つで\R{見}{み}るよ。グラフの\R{傾}{かたむ}きは「\R{変化}{へん|か}の\R{速}{はや}さ」で、2\R{本}{ほん}の\R{線}{せん}が\R{交}{まじ}わる\R{所}{ところ}は、2つの\R{式}{しき}の\R{答}{こた}えだよ。}
\bkPartContents{1}{%% 自動生成（gen.py）
\bktocchap{3-1}{\R{変}{か}わる\R{量}{りょう}と\R{変}{か}わらない\R{量}{りょう}}{ずっと\R{変}{か}わらない\R{数}{かず}と、\R{問題}{もん|だい}ごとに\R{決}{き}まる\R{数}{かず}}{7}{1}{1}{24}{C}
\bktocchap{3-2}{\R{場所}{ば|しょ}を\R{数}{かず}で\R{表}{あらわ}す：\R{座標}{ざ|ひょう}}{「どこ」を2つの\R{数}{かず}で\R{伝}{つた}えよう}{6}{1}{2}{48}{C}
\bktocchap{3-3}{\R{関数}{かん|すう}は「\R{入}{い}れると\R{出}{で}てくる」ルール}{\R{表}{ひょう}・グラフ・\R{式}{しき}は、\R{同}{おな}じことを\R{言}{い}っている}{5}{1}{2}{76}{C}
\bktocchap{3-4}{グラフをかこう・\R{読}{よ}もう}{\R{実験}{じっ|けん}の\R{点}{てん}に、なめらかな\R{線}{せん}を\R{引}{ひ}こう}{7}{1}{1}{104}{C}
\bktocchap{3-5}{グラフの\R{傾}{かたむ}きは「\R{変化}{へん|か}の\R{速}{はや}さ」}{\R{坂}{さか}が\R{急}{きゅう}なほど、\R{速}{はや}く\R{変}{か}わっている}{1}{2}{2}{128}{C}
\bktocchap{3-6}{2\R{本}{ほん}の\R{直線}{ちょく|せん}が\R{交}{まじ}わるところ：\R{連立方程式}{れん|りつ|ほう|てい|しき}}{ふたりが\R{出会}{で|あ}うのは、グラフの\R{交}{まじ}わる\R{点}{てん}}{5}{2}{2}{162}{C}
}%
{\R{実験}{じっ|けん}の\R{表}{ひょう}からグラフをかき、\R{傾}{かたむ}きと\R{式}{しき}を\R{読}{よ}みとれるようになります。「\R{入}{い}れると\R{出}{で}てくる」\R{関数}{かん|すう}の\R{見方}{み|かた}は、\R{第}{だい}5\R{巻}{かん}の\R{微分}{び|ぶん}から\R{第}{だい}9\R{巻}{かん}の\R{制御}{せい|ぎょ}まで、ずっと\R{使}{つか}います。}
\bkpartbegin{1}{\R{関数}{かん|すう}とグラフ}
\bkinputchap{ch3-01}
\bkinputchap{ch3-02}
\bkinputchap{ch3-03}
\bkinputchap{ch3-04}
\bkinputchap{ch3-05}
\bkinputchap{ch3-06}

%% ======================== 第2部 曲がったグラフ ========================
\bkPartOpener{2}{\R{第}{だい}3\R{巻}{かん}\hspace{0.8em}\R{第}{だい}2\R{部}{ぶ}}{\R{曲}{ま}がったグラフ}%
  {172/{おし\R{返}{かえ}す\R{力}{ちから}は、なぜだんだん\R{強}{つよ}くなる？},
   160/{\R{投}{な}げたボールは、なぜこの\R{形}{かたち}に\R{飛}{と}ぶ？},
   148/{ボールが\R{同}{おな}じ\R{高}{たか}さを2\R{回}{かい}\R{通}{とお}るのは、いつ？}}%
  {% 放物線と反比例の曲線（濃い紺の線）
   \begin{scope}[bkOn, line width=1.2pt, line cap=round, line join=round]
     \draw (118,112) -- (172,112);
     \draw[line width=1.6pt, domain=-1:1, samples=40] plot ({145+24*\x}, {112+38*(1-\x*\x)});
     \fill (133,140.5) circle (2.0);% 第13回：ボールの点を放物線の上にぴったり置いた
     \draw[line width=1pt, domain=0.25:1, samples=30, dash pattern=on 2pt off 2pt] plot ({120+12/\x}, {112+9/\x});
   \end{scope}}%
  {まっすぐでないグラフにも、きまりがあるよ。\R{式}{しき}のかっこをはずしたり、\R{形}{かたち}を\R{変}{か}えたりすると、ボールのいちばん\R{高}{たか}い\R{所}{ところ}や、\R{届}{とど}く\R{時刻}{じ|こく}が\R{計算}{けい|さん}で\R{見}{み}つかるよ。}
\bkPartContents{2}{%% 自動生成（gen.py）
\bktocchap{3-7}{\R{反比例}{はん|ぴ|れい}のグラフ}{\R{空気}{くう|き}をおしちぢめるほど、おし\R{返}{かえ}す\R{力}{ちから}は\R{強}{つよ}くなる}{1}{2}{2}{192}{C}
\bktocchap{3-8}{かっこをはずす：\R{展開}{てん|かい}と\R{平方完成}{へい|ほう|かん|せい}}{\R{式}{しき}の\R{形}{かたち}を\R{変}{か}えると、\R{最高点}{さい|こう|てん}が\R{見}{み}えてくる}{1}{2}{3}{218}{C}
\bktocchap{3-9}{\R{放物線}{ほう|ぶつ|せん}と $y=ax^2$}{\R{投}{な}げたボールは、なぜこの\R{形}{かたち}に\R{飛}{と}ぶ？}{1}{2}{3}{256}{C}
\bktocchap{3-10}{ボールはいつその\R{高}{たか}さに\R{届}{とど}く？：2\R{次方程式}{じ|ほう|てい|しき}}{\R{答}{こた}えが2つ\R{出}{で}るのは、\R{上}{のぼ}りと\R{下}{くだ}りがあるから}{1}{2}{3}{290}{C}
}%
{\R{反比例}{はん|ぴ|れい}と\R{放物線}{ほう|ぶつ|せん}のグラフを\R{読}{よ}み、2\R{次方程式}{じ|ほう|てい|しき}を\R{解}{と}けるようになります。\R{式}{しき}の\R{形}{かたち}を\R{変}{か}える\R{力}{ちから}（\R{展開}{てん|かい}・\R{恒等式}{こう|とう|しき}・\R{因数分解}{いん|すう|ぶん|かい}）は、\R{微分}{び|ぶん}・\R{積分}{せき|ぶん}と\R{制御}{せい|ぎょ}の\R{計算}{けい|さん}の\R{土台}{ど|だい}になります。}
\bkpartbegin{2}{\R{曲}{ま}がったグラフ}
\bkinputchap{ch3-07}
\bkinputchap{ch3-08}
\bkinputchap{ch3-09}
\bkinputchap{ch3-10}

%% ======================== 第3部 倍々とけた ========================
\bkPartOpener{3}{\R{第}{だい}3\R{巻}{かん}\hspace{0.8em}\R{第}{だい}3\R{部}{ぶ}}{\R{倍々}{ばい|ばい}とけた}%
  {172/{\R{倍々}{ばい|ばい}でふえると、どれくらい\R{大}{おお}きくなる？},
   160/{\R{半分}{はん|ぶん}、\R{半分}{はん|ぶん}とへっていく\R{量}{りょう}は、いつ0になる？},
   148/{\R{音}{おと}の\R{大}{おお}きさの「デシベル」は、\R{何}{なに}を\R{数}{かぞ}えている？}}%
  {% 倍々の曲線と、けたの目もり（濃い紺の線）
   \begin{scope}[bkOn, line width=1.2pt, line cap=round, line join=round]
     \draw (118,112) -- (172,112);
     \draw (122,108) -- (122,140);% 第12回：3つめの問いの「？」に当たらないよう、たて軸と曲線を140までにした
     \draw[line width=1.6pt, domain=0:46, samples=50] plot ({122+\x}, {112+0.78*2^(\x/9)});
     \foreach \y in {120,130} {\draw[line width=0.9pt] (119,\y) -- (125,\y);}
   \end{scope}}%
  {\R{倍々}{ばい|ばい}でふえる\R{量}{りょう}は、あっという\R{間}{ま}にとても\R{大}{おお}きくなるよ。そんな\R{量}{りょう}は「けたがいくつか」で\R{数}{かぞ}えると、\R{見}{み}やすくなるよ。}
\bkPartContents{3}{%% 自動生成（gen.py）
\bktocchap{3-11}{\R{指数関数}{し|すう|かん|すう}：いまある\R{量}{りょう}に\R{比例}{ひ|れい}してふえる・へる}{\R{倍々}{ばい|ばい}でふえると、あっという\R{間}{ま}に\R{大}{おお}きくなる}{3}{2}{3}{328}{C}
\bktocchap{3-12}{\R{対数}{たい|すう}：けたを\R{数}{かぞ}える\R{道具}{どう|ぐ}}{\R{音}{おと}の\R{大}{おお}きさの「デシベル」は、けたの\R{数}{かぞ}え\R{方}{かた}}{3}{2}{3}{358}{C}
}%
{\R{指数関数}{し|すう|かん|すう}と\R{対数}{たい|すう}を\R{使}{つか}って、ふえる・へる\R{量}{りょう}を\R{計算}{けい|さん}し、けたのちがう\R{量}{りょう}を1つのグラフにかけるようになります。\R{第}{だい}5\R{巻}{かん}の $e$ と、\R{第}{だい}9\R{巻}{かん}の\R{制御}{せい|ぎょ}につながります。}
\bkpartbegin{3}{\R{倍々}{ばい|ばい}とけた}
\bkinputchap{ch3-11}
\bkinputchap{ch3-12}

%% ======================== 第4部 データと統計 ========================
\bkPartOpener{4}{\R{第}{だい}3\R{巻}{かん}\hspace{0.8em}\R{第}{だい}4\R{部}{ぶ}}{データと\R{統計}{とう|けい}}%
  {172/{\R{何回}{なん|かい}もはかると、\R{本当}{ほん|とう}の\R{値}{あたい}に\R{近}{ちか}づける？},
   160/{ずれた\R{点}{てん}に、どの\R{直線}{ちょく|せん}を\R{引}{ひ}けばいい？},
   148/{サイコロをたくさん\R{足}{た}すと、どんな\R{形}{かたち}になる？}}%
  {% ヒストグラムとつりがね形（濃い紺の線）
   \begin{scope}[bkOn, line width=1.1pt, line join=round]
     \draw (118,112) -- (172,112);
     \foreach \i/\h in {0/4,1/10,2/20,3/30,4/34,5/30,6/20,7/10,8/4} {\draw (121+\i*5.6,112) rectangle ++(5.6,\h);}
     \draw[line width=1.6pt, domain=-4.5:4.5, samples=60] plot ({146.2+5.6*\x}, {112+36*exp(-\x*\x/4.4)});
   \end{scope}}%
  {はかった\R{数}{かず}はいつも\R{少}{すこ}しずつずれるよ。そのずれを\R{表}{ひょう}とグラフにまとめ、「だいたいこのくらいずれる」を1つの\R{数}{かず}で\R{言}{い}う\R{方法}{ほう|ほう}を\R{学}{まな}ぶよ。}
\bkPartContents{4}{%% 自動生成（gen.py）
\bktocchap{3-13}{データのちらばり：ヒストグラムと\R{平均}{へい|きん}}{\R{何回}{なん|かい}もはかって、\R{本当}{ほん|とう}の\R{値}{あたい}に\R{近}{ちか}づこう}{7}{2}{2}{388}{C}
\bktocchap{3-14}{ばらつきを1つの\R{数}{かず}で：\R{標準偏差}{ひょう|じゅん|へん|さ}}{「だいたいこのくらいずれる」を\R{数}{かず}で\R{言}{い}う}{7}{3}{3}{416}{C}
\bktocchap{3-15}{2つの\R{量}{りょう}の\R{関係}{かん|けい}：\R{散布図}{さん|ぷ|ず}と\R{相関}{そう|かん}}{\R{背}{せ}の\R{高}{たか}い\R{人}{ひと}は、\R{足}{あし}も\R{速}{はや}い？}{7}{3}{3}{442}{C}
\bktocchap{3-16}{データに\R{直線}{ちょく|せん}を\R{引}{ひ}く：\R{最小二乗法}{さい|しょう|に|じょう|ほう}}{ずれがいちばん\R{小}{ちい}さくなる\R{線}{せん}は、どれ？}{7}{3}{3}{474}{C}
\bktocchap{3-17}{\R{偶然}{ぐう|ぜん}のばらつき：\R{確率}{かく|りつ}と\R{正規分布}{せい|き|ぶん|ぷ}}{サイコロをたくさん\R{足}{た}すと、つりがねの\R{形}{かたち}に}{7}{3}{3}{498}{C}
}%
{はかったデータから、\R{平均}{へい|きん}とばらつきを\R{計算}{けい|さん}し、\R{点}{てん}にいちばん\R{合}{あ}う\R{直線}{ちょく|せん}を\R{引}{ひ}けるようになります。\R{計測}{けい|そく}と\R{品質}{ひん|しつ}の\R{管理}{かん|り}、\R{第}{だい}9\R{巻}{かん}の\R{信号}{しん|ごう}の\R{処理}{しょ|り}につながります。}
\bkpartbegin{4}{データと\R{統計}{とう|けい}}
\bkinputchap{ch3-13}
\bkinputchap{ch3-14}
\bkinputchap{ch3-15}
\bkinputchap{ch3-16}
\bkinputchap{ch3-17}

%% ======================== 巻末 ========================
\bkinputchap{ch3-back}
